\subsection{Integer-valued polynomials and the small-prime estimate}
For a prime $p$ and integer $D\ge1$, write
\[
 \ell_p(D)=\max\{e\in\N:p^e\le D\}=\floor{\log_p D}.
\]
The integer characterization is preferable in a formal development.

\begin{lemma}[Losses for integer-valued polynomials]
\label{lem:integervalued}
Let $A\in\Qp[x]$ have degree at most $d$ and map $\Zp$ into $\Zp$.
Then its coefficients in the basis $\binom{x}{k}$ are integral and
\[
 v_p(A(x)-A(y))\ge v_p(x-y)-\ell_p(\max(1,d))\qquad(x,y\in\Zp).
\]
Its $r$th derivative is integer-valued up to a loss of at most
$r\ell_p(\max(1,d))$. If $P(\Zp)\subseteq p^{-A_0}\Zp$ and
$\deg P\le d$, then
\begin{equation}\label{eq:tauivbound}
 v_p(\tau(P))\ge-A_0-4\ell_p(d+1)-v_p(24).
\end{equation}
\gid{Zeta5.Local.integer_valued_estimates}
\end{lemma}
\begin{proof}
The binomial coefficients of $A$ are its finite forward differences at
zero, hence integral. Vandermonde's identity, together with
$\binom{z}{j}=z\binom{z-1}{j-1}/j$, gives the difference estimate.
For derivatives expand $\binom{x+z}{k}$ by Vandermonde. The coefficient
of $z^r$ in $\binom zj$ is a sum of products of $r$ reciprocals of
positive integers at most $j$: use
$\binom zj=(-1)^{j-1}(z/j)\prod_{u=1}^{j-1}(1-z/u)$.
Each product has valuation at least $-r\ell_p(\max(1,d))$; multiplication
by $r!$ cannot worsen it.

The formal generating identity
$L((1+z)^x)=\log(1+z)/z$ gives
$L(\binom{x}{k})=(-1)^k/(k+1)$. Apply the third-derivative estimate,
expand the resulting polynomial in the binomial basis, and apply this
last identity. A further loss at most $\ell_p(d+1)$ and the factor $24$
give \eqref{eq:tauivbound}. Scaling by $p^{A_0}$ treats the stated range.
\end{proof}

\begin{proposition}[Small-prime rational estimate; source Lemma 3.3]
\label{prop:smallprime}
For any prime $p$, let $A\in\Qp[x]$ be integer-valued on $\Zp$ with
$\deg A\le d$, and let $K\ge1$. Set
\[
 g(x)=\frac{(K!)^2A(x)}{\prod_{-K\le r\le K,\ r\ne0}(x-r)}.
\]
Then
\begin{equation}\label{eq:smallprime}
 v_p^{\mathrm G}(\tau_X(g))
 \ge-6\ell_p(\max(2K,d+1))-v_p(24).
\end{equation}
\gid{Zeta5.Local.small_prime_functional_bound}
\end{proposition}
\begin{proof}
Let $S=\{-K,\ldots,-1,1,\ldots,K\}$ and $L_0=\ell_p(2K)$.
Write $g=P+\sum_{r\in S}c_r/(x-r)$. The residue is, up to sign,
\[
 c_r=\frac{(K!)^2rA(r)}{(K+r)!(K-r)!},\qquad v_p(c_r)\ge-L_0.
\]
For each power $p^a\le2K$ the difference between the two denominator
factorial floor sums and $2\floor{K/p^a}$ is at most one. Larger powers
contribute zero, which proves the bound.

Suppose first $x$ lies outside every ball $r+p^{L_0+1}\Zp$. At level
$p^a$, $a\le L_0$, each of the two length-$K$ integer intervals in $S$
contains at most $\floor{K/p^a}+1$ elements congruent to $x$. There are
no roots with $v_p(x-r)>L_0$. Counting by levels therefore gives
$v_p(g(x))\ge-2L_0$. Each principal part has the same lower bound, so
$v_p(P(x))\ge-2L_0$.

Now suppose $x\in r+p^{L_0+1}\Zp$. This center is unique, since distinct
roots differ by an integer of absolute value at most $2K$. Put
\[
 F_r(x)=\frac{(K!)^2}{\prod_{s\in S\setminus\{r\}}(x-s)}.
\]
For every $s\ne r$, $v_p(x-s)=v_p(r-s)\le L_0$. Thus
$v_p(F_r(x))=v_p(F_r(r))\ge-L_0$. If $q=v_p(x-r)>L_0$, expand the
finite product
\[
 \frac{F_r(x)}{F_r(r)}
 =\prod_{s\ne r}\left(1+\frac{x-r}{r-s}\right)^{-1}.
\]
Every factor is $1$ modulo $p^{q-L_0}$, whence
$v_p(F_r(x)-F_r(r))\ge q-2L_0$. Subtracting the principal part at $r$
yields
\[
 \frac{A(x)F_r(x)-A(r)F_r(r)}{x-r}
 =\frac{(A(x)-A(r))F_r(x)}{x-r}
  +\frac{A(r)(F_r(x)-F_r(r))}{x-r}.
\]
Lemma~\ref{lem:integervalued} bounds its valuation below by
$-L_0-\max(L_0,\ell_p(\max(1,d)))$. All other principal parts are bounded
below by $-2L_0$. Continuity of the polynomial $P$ includes the centers.
Consequently
\[
 P(\Zp)\subseteq p^{-L_0-M_0}\Zp,
 \qquad M_0=\max(L_0,\ell_p(\max(1,d))).
\]
Here $M_0$ is local to this proof, not the real-potential constant used
later. Since $\deg P\le d$, \eqref{eq:tauivbound} gives the required
$-6\ell_p(\max(2K,d+1))-v_p(24)$ bound for its value. For a principal
part, $d(r)\le K$ implies
$v_p(\Hfive{d(r)})\ge-5\ell_p(K)$, and its coefficient of $X$ has
valuation at least $-L_0$. This gives the same bound coefficientwise.
The ultrametric inequality finishes the proof.
\end{proof}

Define
\[
 q_0(t)=1,\qquad q_i(t)=\frac{(-1)^i2tD_{i-1}(t)}{(2i)!}\ (i\ge1),
 \qquad f_i(t)=\left(\frac{D_N(t)}{(N!)^2}\right)^3q_i(t).
\]
The identities
\[
 q_i(-x^2)=\binom{x+i}{2i}+\binom{x+i-1}{2i},\qquad
 \frac{D_N(-x^2)}{(N!)^2}=\binom{N-x}{N}\binom{N+x}{N}
\]
show that $f_i(-x^2)$ is integer-valued on every $\Zp$. The triangular
change from monomials to $q_i$ gives exactly
\begin{equation}\label{eq:Fbasis}
 F_K(X)=\det\left[(K!)^2\mu_X\left(\frac{f_i(t)f_j(t)}{D_K(t)}\right)
              \right]_{0\le i,j<h}.
\end{equation}
The pullback numerator $x^5f_i(-x^2)f_j(-x^2)$ has degree at most
$12N+4h+1$, so $\max(2K,d+1)\le5K$. Applying
Proposition~\ref{prop:smallprime} entrywise and expanding the determinant
proves the universal bound
\begin{equation}\label{eq:Fsmall}
 v_p^{\mathrm G}(F_K)\ge-6h\ell_p(5K)-h v_p(24)
 \qquad\text{for every prime }p.
\end{equation}
\gid{Zeta5.small_prime_determinant_bound}

\section{Local determinant estimates}
\label{sec:detlocal}
Fix an integer $M\ge40$ and assume $K\in40\Z_{>0}$ and $K\ge200M^2$.
The small-prime estimate is used for $p\le K/M$. Two separate bases give
stronger estimates above this cutoff.

\subsection{The inner range: allocation, degree checks, and every weight case}
Suppose
\begin{equation}\label{eq:innerhyp}
 K/M<p\le K/3.
\end{equation}
Then $p>200M$, $p^2>5K$, and $p\ge7$. Put $m=(p-1)/2$, and for
$1\le a\le m$ define
\[
 \ell_A(a)=\#\{1\le j\le A:j\equiv a\text{ or }-a\pmod p\},
 \qquad m_A=\floor{A/p}.
\]
If a row polynomial has a zero of order $\nu_i(a)$ at $t=-a^2$, the
residue substitutions $x=c+pz$ in \eqref{eq:pullback} and
\eqref{eq:distribution} give the following candidate entry bounds:
\begin{align}
 \nu_i(a)+\nu_j(a)+6\ell_N(a)-\ell_K(a)-4
 &\quad (c=a\text{ or }p-a),\label{eq:ordinarysource}\\
 2\nu_i(0)+2\nu_j(0)+12m_N-2m_K+1
 &\quad (c=0).\label{eq:zerosource}
\end{align}
We justify the degree and completed-domain conditions after specifying
the basis; the displayed powers alone are not yet the full argument.

Set $L_0=4M+10$ and $b_a=3\ell_N(a)$. There are unique integers $T,E$
with
\begin{equation}\label{eq:allocation}
 mT+E=h-L_0+3(N-m_N),\qquad0\le E<m.
\end{equation}
Order the ordinary classes by nonincreasing $\ell_K(a)$, breaking ties
by increasing $a$, and give $\epsilon_a=1$ to the first $E$ classes and
$0$ to the others. Define
\[
 L_a=T-b_a+\epsilon_a,\qquad Z_a=L_a+b_a=T+\epsilon_a.
\]
Because $\sum_{a=1}^m\ell_N(a)=N-m_N$,
$L_0+\sum_{a=1}^mL_a=h$. With $x=K/p\in[3,M)$, \eqref{eq:allocation}
implies
\[
 2Hx-\frac{21}{20}<T<2Hx,\qquad b_a\le6\alpha x+3.
\]
For the lower bound, write the numerator in \eqref{eq:allocation} as
$HK-L_0-3m_N$, use $2E/(p-1)<1$, and use
$(2L_0+6m_N)/(p-1)<1/20$. For the upper bound,
$2L_0+6m_N>2Hx$. These inequalities follow from $p>200M$,
$m_N\le\alpha M$, and $L_0=4M+10$. Consequently
\[
 T-b_a>2\lambda x-\frac{81}{20}\ge\frac32,
\]
so each ordinary $L_a$ is a positive integer.

For $0\le a\le m$ and $0\le i<L_a$, let
\begin{equation}\label{eq:innerbasis}
 E_{a,i}(t)=\prod_{\substack{0\le c\le m\\c\ne a}}(t+c^2)^{L_c}
                      (t+a^2)^i.
\end{equation}
All degrees are less than $h$. The factors $t+c^2$ are pairwise coprime
modulo $p$. In the Chinese remainder decomposition for their powers,
the local blocks of \eqref{eq:innerbasis} are triangular with unit
diagonal. Thus these $h$ polynomials form a $\Zp$-unimodular basis of
polynomials of degree less than $h$.

Define half-integer weights
\begin{align}
 w_{a,i}&=i+b_a-\frac{\ell_K(a)+4}{2} &&(a>0),\label{eq:innerweights}\\
 w_{0,i}&=\min\left(2i+6m_N-m_K+\frac12,
                   \min_{c>0}\left(Z_c-\frac{\ell_K(c)+4}{2}\right)\right),
 \label{eq:zeroweights}\\
 \gamma_p^{\mathrm{in}}&=2\sum_{a=0}^m\sum_{i=0}^{L_a-1}w_{a,i}\in\Z.
 \label{eq:gammain}
\end{align}

\begin{proposition}[Inner valuation; source Proposition 4.1]
\label{prop:inner}
Under \eqref{eq:innerhyp},
$v_p^{\mathrm G}(\Delta_K)\ge\gamma_p^{\mathrm{in}}$.
\gid{Zeta5.inner_determinant_valuation}
\end{proposition}
\begin{proof}
\emph{The local summands really have the required integral form.}
After $x=c+pz$, a near denominator factor is $p$ times a linear factor
with an integer zero. The near integers have absolute value at most
$M+1$ and mutual nonzero differences of absolute value less than $p$,
so those differences are units. Their $d(r)$ indices are less than $p$.
A far factor is a unit times a geometric series in $pz$. After extracting
the indicated power of $p$, the ordinary summand has the form
$\sum_{j\ge0}p^jU_j/T$ of Proposition~\ref{prop:integralextension}.
Its leading numerator has degree at most
\[
 \nu_i(a)+\nu_j(a)+6\ell_N(a)\le2(T+1)<\frac{23}{5}M+2<p+1.
\]
Here every row satisfies $\nu_i(a)+b_a\le Z_a\le T+1$.
At the zero source, $x^5$ contributes $p^5z^5$, a zero of order $\nu$
in $t$ contributes $p^{2\nu}$, and the $\pm jp$ poles contribute the
$2m_K$ denominator factors. Its leading numerator degree is at most
\[
 5+2\nu_i(0)+2\nu_j(0)+12m_N
 \le5+4L_0+12m_N\le\frac{169}{10}M+45<p+1.
\]
The outside distribution factor is $p^{-4}$. Together with
$Y\in\Zp[X]$, these observations prove
\eqref{eq:ordinarysource}--\eqref{eq:zerosource} for the \emph{whole}
summands, including their analytic polynomial quotients.

\emph{Each source dominates the assigned weight for each row.}
The following table lists the half-bound contributed by one row. Adding
the two row entries in the same source column gives the entry bound.
\begin{center}
\begin{tabular}{lll}
\toprule
Row & Ordinary source $c>0$ & Zero source\\
\midrule
$(a,i)$, $a=c>0$ & $w_{a,i}$ & $2L_0+6m_N-m_K+1/2$\\
$(a,i)$, $a\ne c$, $a>0$ & $Z_c-(\ell_K(c)+4)/2$
                       & $2L_0+6m_N-m_K+1/2$\\
$(0,i)$ & $Z_c-(\ell_K(c)+4)/2$ & $2i+6m_N-m_K+1/2$\\
\bottomrule
\end{tabular}
\end{center}
The zero row is covered directly by the minimum in
\eqref{eq:zeroweights}. For an ordinary row $a\ne c$, since
$i\le L_a-1$, the difference between the table entry and its assigned
weight is at least
\[
 Z_c-Z_a+1+\frac{\ell_K(a)-\ell_K(c)}2.
\]
The counts $\ell_K$ take two consecutive values. The exhaustive extra-row
cases are
\begin{center}
\begin{tabular}{ccl}
\toprule
$\epsilon_a$ & $\epsilon_c$ & Lower bound for the difference\\
\midrule
0 & 0 & $1/2$\\
1 & 1 & $1/2$\\
0 & 1 & $3/2$\\
1 & 0 & $0$, because the allocation order gives $\ell_K(a)\ge\ell_K(c)$\\
\bottomrule
\end{tabular}
\end{center}
At the zero source an ordinary row has half-bound at least
$2L_0-m_K+1/2\ge7M+41/2$. On the other hand
$w_{a,i}<T+1\le2HM+1$. The zero source therefore also dominates.
These checks exhaust every row and every source.

Consequently every matrix entry in this basis has Gauss valuation at
least the sum of its two assigned row weights. In a determinant term
those weights occur twice in total, so its valuation is at least
$\gamma_p^{\mathrm{in}}$. The ultrametric inequality gives the determinant
bound. A unimodular basis change multiplies its determinant by a square
of a unit and does not alter the valuation.
\end{proof}

\subsection{The outer range: polynomial correction and its actual rank}
Assume for this subsection that
\begin{equation}\label{eq:outerhyp}
 p\ge7,\quad p\le K<3p,\quad p^2>2K,\quad
 2N<p,\quad5N\le2p-2.
\end{equation}
These conditions hold for $K/3<p\le K$ under the global admissibility
